\documentclass{article}

\usepackage{amsmath,amssymb}
\usepackage{xurl}
\usepackage[colorlinks=true,linkcolor=blue,citecolor=blue,urlcolor=blue]{hyperref}
\setlength{\emergencystretch}{3em}

% Shared commands for notes in docs/paper-gaps/.

\DeclareUnicodeCharacter{2080}{\ifmmode{}_{0}\else\textsubscript{0}\fi}
\DeclareUnicodeCharacter{2081}{\ifmmode{}_{1}\else\textsubscript{1}\fi}
\DeclareUnicodeCharacter{2082}{\ifmmode{}_{2}\else\textsubscript{2}\fi}
\DeclareUnicodeCharacter{2099}{\ifmmode{}_{n}\else\textsubscript{n}\fi}

\newcommand{\Lean}{\textsc{Lean}}
\ExplSyntaxOn
\NewDocumentCommand{\leanid}{m}{
  \ifmmode
    \textsf{\detokenize{#1}}
  \else
    \group_begin:
    \urlstyle{sf}
    \tl_set:Nn \l_tmpa_tl {#1}
    \tl_replace_all:Nnn \l_tmpa_tl {\_} {_}
    \exp_args:NV \path \l_tmpa_tl
    \group_end:
  \fi
}
\ExplSyntaxOff
\providecommand{\tr}{\operatorname{tr}}
\newcommand{\tnleanIssueBase}{https://github.com/LionSR/TNLean/issues/}
\newcommand{\tnleanPrBase}{https://github.com/LionSR/TNLean/pull/}
\newcommand{\tnleanDocBase}{https://sirui-lu.com/TNLean/docs/}
\newcommand{\ghissue}[1]{\href{\tnleanIssueBase#1}{GitHub issue \##1}}
\newcommand{\ghpr}[1]{\href{\tnleanPrBase#1}{PR \##1}}
\newcommand{\doclink}[1]{\href{\tnleanDocBase#1.html}{\path{#1}}}

% Dirac notation and the identity operator, as in blueprint/src/macros/common.tex.
% \providecommand keeps note-local definitions authoritative.
\usepackage{dsfont}
\providecommand{\ket}[1]{|#1\rangle}
\providecommand{\bra}[1]{\langle #1|}
\providecommand{\braket}[2]{\langle #1 | #2 \rangle}
\providecommand{\ketbra}[2]{|#1\rangle\!\langle #2|}
\providecommand{\Id}{\mathds{1}}
\providecommand{\one}{\mathds{1}}
\providecommand{\C}{\mathbb{C}}
\providecommand{\MN}[1]{M_{#1}(\C)}
\providecommand{\MD}{M_D(\C)}

% External referencing for paper sources.  The build helper writes sanitized
% auxiliary files under build/paper-aux/ and excludes duplicated source labels.
\usepackage{xr}

% Import a generated paper auxiliary file with a caller-supplied label prefix.
% The candidate paths support builds from docs/paper-gaps/, the repository root,
% and build/paper-aux/ itself.
\newcommand{\importpaperaux}[2]{%
  \IfFileExists{../../build/paper-aux/#2.aux}{%
    \externaldocument[#1]{../../build/paper-aux/#2}%
  }{%
    \IfFileExists{build/paper-aux/#2.aux}{%
      \externaldocument[#1]{build/paper-aux/#2}%
    }{%
      \IfFileExists{#2.aux}{%
        \externaldocument[#1]{#2}%
      }{}%
    }%
  }%
}

% General external reference macros
\newcommand{\paperref}[2]{\ref{#1-#2}}
\newcommand{\papereqref}[2]{(\ref{#1-#2})}

% CPSV16 source (1606.00608 / MPDO-22-12-17-2.tex) external document and shortcuts
\importpaperaux{cpsv16-}{1606.00608/MPDO-22-12-17-2-tnlean}
\newcommand{\cpsvref}[1]{\paperref{cpsv16}{#1}}
\newcommand{\cpsveqref}[1]{\papereqref{cpsv16}{#1}}

% Verdict marker: \gapnote{<kind>}{<status>}. Kind is the policy
% classification (clarification, local-correction, scope-restriction,
% unfaithful, false-source, open-gap); status is open, wip (elimination
% underway), resolved, or historical. Machine-read by the site index;
% prints a verdict line.
\newcommand{\gapnote}[2]{\par\noindent\textbf{Verdict:} #1 (#2).\par}


\title{The CPSV16 Unit-Weight Convention and Renormalization-Fixed-Point Scale}
\author{}
\date{2026-08-15}

\begin{document}
\maketitle

\section{Two scale conventions in the source}

The canonical-form construction of
Cirac--P\'erez-Garc\'ia--Schuch--Verstraete first writes a tensor as
\begin{align}
    A^i
    &=\bigoplus_k \mu_k A_k^i,
\end{align}
where every $A_k$ is normal and its transfer map has spectral radius one
\cite[lines~214--246]{Cirac2016MPDO_arXiv}. Since the states are not normalized
there, line~246 adopts the standing convention
\begin{align}
    |\mu_k|&\leq 1 &&\text{for every $k$},
    & |\mu_{k_0}|&=1 &&\text{for some $k_0$}.
    \label{eq:unit-weight}
\end{align}
The formal BNT predicate records the second clause through
\leanid{MPSTensor.IsBNTCanonicalForm.weight_unit_exists}. This is one global
unit-weight witness, not one witness in every sector; the latter distinction is
discussed separately in
\path{docs/paper-gaps/cpsv16_global_vs_persector_unit_witness.tex}.

Definition~4.1 fixes the scale of the tensor in a different way. For a tensor
$M$ already in canonical form and generating MPDOs, it requires
trace-preserving completely positive maps relating the one-site and two-site
physical closures \cite[lines~645--660]{Cirac2016MPDO_arXiv}. Trace preservation
gives
\begin{align}
    \mathcal T_M^2&=\mathcal T_M,
    \label{eq:physical-idempotent}
\end{align}
where $\mathcal T_M=\sum_i M^{ii}$ is the physical-trace transfer; this is the
calculation at lines~1333--1340. The formal statement is
\leanid{MPOTensor.physTraceTransfer_sq_of_isRFPViaTS}.

These requirements concern two transfer objects with different scaling laws.
For the doubled-index MPS transfer map
\begin{align}
    \mathcal E_M(X)
    &=\sum_i M^iX(M^i)^\dagger,
\end{align}
one has
\begin{align}
    \mathcal E_{cM}
    &=|c|^2\mathcal E_M.
    \label{eq:mps-transfer-scaling}
\end{align}
For the physical-trace transfer, linearity gives
\begin{align}
    \mathcal T_{cM}
    &=c\mathcal T_M.
    \label{eq:physical-transfer-scaling}
\end{align}
Equation \eqref{eq:unit-weight} normalizes the first object by choosing normal
BNT representatives and placing the overall scale in their coefficients.
Equation \eqref{eq:physical-idempotent} constrains the second object at the
scale of the displayed tensor. The related up-to-scalar treatment of source
zero correlation length is recorded in
\path{docs/paper-gaps/cpsv16_zcl_canonical_form_normalization.tex}.

\section{Fixed-representative scale pinning}

By \eqref{eq:physical-transfer-scaling}, replacing $M$ by $cM$ replaces
$\mathcal T_M$ by $c\mathcal T_M$. If
$\mathcal T_M^2=\mathcal T_M\ne0$ and the nonzero rescaling $cM$ also satisfies
the same literal idempotence equation, then
\begin{align}
    c^2\mathcal T_M
    &=(c\mathcal T_M)^2
      =c\mathcal T_M,
\end{align}
so $c=1$. Thus Definition~4.1 is not invariant under an arbitrary nonzero
rescaling of a fixed representative.

By contrast, the simplicity definition at lines~815--822 is phrased through
the elements of a BNT: the tensor is simple when none of those elements has
nilpotent ket-against-bra contraction. Multiplying one BNT element by a nonzero
scalar does not change nilpotency. This observation concerns the nilpotency
clause alone; it does not imply that an arbitrary complex rescaling preserves
the MPDO condition or either full simplicity predicate.

Indeed, if $\rho_N(M)$ denotes the closed length-$N$ MPO, then the exact law is
\begin{align}
    \rho_N(cM)&=c^N\rho_N(M),
    \label{eq:closed-mpo-scaling}
\end{align}
not $|c|^{2N}\rho_N(M)$. Consequently positivity is preserved in both
directions by a strictly positive real rescaling $rM$, and the formal theorems
\leanid{MPOTensor.isMPDO_smul_ofReal_iff} and
\leanid{MPOTensor.isSourceSimple_smul_ofReal_iff} establish the corresponding
invariances. They do not assert invariance under arbitrary nonzero complex
scalars, and they do not imply that the rescaled tensor satisfies the same
literal fixed-scale RFP equations of Definition~4.1.

The declaration \leanid{MPOTensor.IsSourceSimple} records the
normalization-free source condition: it requires every positive-length
generated MPO to be nonzero and asks for some positive blocking and a basis of
normal tensors whose physical-trace transfers are all non-nilpotent. It does
not impose the line-246 unit-weight convention.

By contrast, \leanid{MPOTensor.IsSimpleCanonicalForm} and
\leanid{MPOTensor.IsSimple} are normalized fixed-representative interfaces. The
first is a predicate on one displayed canonical-form representative. The second
permits positive physical blocking, but still asks the resulting blocked tensor
itself to admit such a normalized witness. Neither declaration is a quotient by
nonzero scalar rescaling, and no equivalence with
\leanid{MPOTensor.IsSourceSimple} is asserted. There is, however, a sanctioned
one-way bridge: \leanid{MPOTensor.IsSimple.isSourceSimple} derives source
simplicity from normalized simplicity under the explicit hypothesis
\begin{align}
    \rho_N(M)&\ne0 &&\text{for every }N>0.
    \label{eq:positive-length-nonzero}
\end{align}
The extra hypothesis cannot be omitted from the formal statement because the
normalized sector weights may cancel at particular lengths. For a fixed tensor,
a theorem may also quantify over every normalized BNT witness rather than rely
on one chosen presentation; the dimer obstruction below has precisely this
form.

\section{The displayed dimer representative}

The project tensor $R$ gives a concrete instance of this tension. Its
doubled-index MPS transfer map satisfies
\begin{align}
    \mathcal E_R^2
    &=\frac{337}{512}\mathcal E_R.
\end{align}
Consequently the displayed one-block CPSV canonical form is
\begin{align}
    R^{\bullet\bullet}
    &=\mu\widehat R,
    & \mu&=\sqrt{\frac{337}{512}},
\end{align}
where the transfer map of the retained block $\widehat R$ is idempotent and
hence has spectral radius one. This presentation has $0<\mu<1$, so it does not
itself supply the global unit-weight witness in \eqref{eq:unit-weight}.

On the other hand, the physical-trace transfer of the displayed tensor $R$ is
a nonzero idempotent: its $(0,0)$ entry is $1/2$. The retained block has
\begin{align}
    \mathcal T_{\widehat R}
    &=\mu^{-1}\mathcal T_R.
\end{align}
Since $\mu\ne1$, the scale-pinning calculation above shows that this normalized
representative does not satisfy the same literal physical-trace idempotence
equation as $R$.

The obstruction is stronger than this displayed-presentation mismatch. For
every positive blocking length $L$, the doubled-index transfer map of the
blocked tensor satisfies
\begin{align}
    \mathcal E_{R^{[L]}}^2
    &=\left(\frac{337}{512}\right)^L\mathcal E_{R^{[L]}}.
    \label{eq:blocked-quasi-idempotence}
\end{align}
Suppose that $R^{[L]}$ admitted any normalized horizontal BNT witness. The
gauge to that witness transports \eqref{eq:blocked-quasi-idempotence}. The
global unit-weight condition from line~246 supplies a copy whose weight has
modulus one. Restricting the transported equation to that copy removes the
weight. Its normal representative is left canonical, so its transfer map is
trace preserving. Taking the trace of the quasi-idempotence equation at the
identity forces
\begin{align}
    \left(\frac{337}{512}\right)^L&=1.
\end{align}
This contradicts $0<(337/512)^L<1$. Hence no positive blocking of $R$ admits a
normalized horizontal canonical form. Since the formal simple-canonical-form
predicate includes such a horizontal witness, the theorem
\leanid{R_not_isSimple} in the namespace
\leanid{MPOTensor.RescalingStableLengthDependentRFP} proves
\begin{align}
    \lnot\,\leanid{MPOTensor.IsSimple}\,R.
\end{align}
The argument quantifies over the arbitrary normalized BNT witness allowed by
the formal predicate; it does not depend on the displayed one-block
presentation.

This failure is not a nilpotency result. The retained-block theorem proves
that the displayed canonical presentation has non-nilpotent physical-trace
transfer. Together with positivity and the active one-site BNT refinement, it
witnesses
\begin{align}
    \leanid{MPOTensor.IsSourceSimple}\,R.
\end{align}
The source quantifier is existential and is witnessed by $L=1$; no assertion is
made about every positive blocking. The normalized obstruction instead combines
quasi-idempotence of the doubled-index transfer, the line-246 unit-weight copy,
and trace preservation of its normal representative. It uses none of the
vertical coefficient identities discussed below.

\section{The vertical one-label identity}

The same scale bookkeeping is visible in the algebraic form of
Theorem~4.14. Its length-one condition is
\begin{align}
    m_\gamma
    &=\sum_{\alpha,\beta}
      c^{(1)}_{\alpha,\beta,\gamma}m_\alpha m_\beta,
\end{align}
where $m_\alpha=\tr(\mu_\alpha)$
\cite[lines~972--993]{Cirac2016MPDO_arXiv}. In the displayed one-label dimer
presentation,
\begin{align}
    c^{(1)}_{0,0,0}&=\frac{32}{25},
    &m_0&=\frac{25}{32},
\end{align}
so the nonzero solution is fixed by $(c^{(1)}_{0,0,0})^{-1}$ rather than by a
unit-weight convention. This is a statement about that explicit presentation,
not an invariance theorem for coefficients under changes of presentation.

\section{Formalization boundary}

\textbf{Verdict: two distinct formal interfaces.} The source-facing
predicate \leanid{MPOTensor.IsSourceSimple} is true for $R$, with positive
blocking length $L=1$. The normalized fixed-representative predicate
\leanid{MPOTensor.IsSimple} is false for $R$: every positive blocking has
quasi-idempotence scalar $\bigl(337/512\bigr)^L<1$, while every normalized horizontal BNT
witness would force that scalar to be one. The latter conclusion is independent
of the choice of normalized BNT witness.

The normalized verdict remains representative-dependent. The line-246 field
should not be deleted: it is the source's standing canonical-form convention.
The obstruction does not produce a nilpotent BNT block, does not use the
vertical coefficient family, and does not refute source simplicity. Conversely,
the one-site source-simplicity witness does not assert that every positive
blocking is source-simple, that arbitrary complex rescaling preserves the MPDO
condition, that positive rescaling preserves the same fixed-scale RFP equations,
or that the two simplicity interfaces are equivalent. The valid comparison is
the conditional bridge \eqref{eq:positive-length-nonzero}, together with
strictly positive real rescaling invariance of the MPDO and source-simple
predicates.

The declaration-free compatibility shell
\path{TNLean/MPS/MPDO/RescalingStableLengthDependentRFPBNTClause.lean} has been
deleted. The live explicit BNT declarations remain in
\path{TNLean/MPS/MPDO/RescalingStableExplicitVerticalBNT.lean}.

\bibliographystyle{alpha}
\bibliography{references}

\end{document}
